% Gerado por build-tex.py a partir de apostila/indice-figuras.md.
\begingroup\small
\begin{longtable}{@{}r>{\raggedright\arraybackslash}p{0.68\linewidth}r@{}}
\toprule
Figura & Legenda & Página \\
\midrule
\endhead
\bottomrule
\endfoot
1 & \url{https://embarcados.com.br/wp-content/uploads/2016/08/Acabamento-de-} superf\%C3\%ADcie-destaque-1.jpg.webp & \pageref{fig:01} \\
2 & \url{https://resources.altium.com/sites/default/files/inline-images/pcb-silk-3.png} & \pageref{fig:02} \\
3 & Hot Air Solder Level-A & \pageref{fig:03} \\
4 & Hot Air Solder Level-B & \pageref{fig:04} \\
5 & Imersão em Estanho-A & \pageref{fig:05} \\
6 & Imersão em Estanho-B & \pageref{fig:06} \\
7 & Imersão em Prata-A & \pageref{fig:07} \\
8 & Imersão em Prata-B & \pageref{fig:08} \\
9 & Organic Solderability Preservative-A & \pageref{fig:09} \\
10 & Organic Solderability Preservative-B & \pageref{fig:10} \\
11 & Electroless Nickel Immersion Gold) - A & \pageref{fig:11} \\
12 & Electroless Nickel Immersion Gold) - B & \pageref{fig:12} \\
13 & Electroless Nickel Electroless Palladium Immersion Gold-A & \pageref{fig:13} \\
14 & Electroless Nickel Electroless Palladium Immersion Gold-B & \pageref{fig:14} \\
15 & Hard Gold-A & \pageref{fig:15} \\
16 & Hard Gold-B & \pageref{fig:16} \\
17 & Separação dos componentes em grupos de função & \pageref{fig:17} \\
18 & Plano de terra e furos de passagem & \pageref{fig:18} \\
19 & Cada seção do circuito deve ter seu plano de terra, havendo a interligação entre eles & \pageref{fig:19} \\
20 & A ligação de todos os grupos na mesma linha de alimentação e terra não é recomendada & \pageref{fig:20} \\
21 & É recomendado o uso de interfaces entre seções do sistema & \pageref{fig:21} \\
22 & Gerenciador de Projetos. Fonte: \url{https://docs.kicad.org/8.0/en/kicad/kicad.html#:~:text=Usando%20o} \%20gerenciador,editores\%20e\%20ferramentas & \pageref{fig:22} \\
23 & Workflow básico KiCAD. fonte: \url{https://www.slideshare.net/baoshi1/why-} and-how-to-switch-to-kicad & \pageref{fig:23} \\
24 & Exemplo de símbolo & \pageref{fig:24} \\
25 & Exemplo de footprint & \pageref{fig:25} \\
26 & Exemplo de modelos 3D & \pageref{fig:26} \\
27 & Exemplo de indentificação do Part na JLCPCB & \pageref{fig:27} \\
28 & Identificação do componente a ser utilizado na lista BOM para fabricação na JLCPCB & \pageref{fig:28} \\
29 & Configuração das bibliotecas de símbolos & \pageref{fig:29} \\
30 & Interface de configuração das restrições de desenho da PCB & \pageref{fig:30} \\
31 & A64-OLinuXino é um computador de placa única que roda Linux e Android & \pageref{fig:31} \\
32 & Uma câmera infravermelha de imagem térmica de código aberto & \pageref{fig:32} \\
33 & ANAVI Flex é um Raspberry Pi Shield de código aberto para prototipagem de IoT e aplicações de automação residencial & \pageref{fig:33} \\
34 & ANAVI Infrared pHAT é uma placa adicional que converte o Raspberry Pi em um poderoso controle remoto & \pageref{fig:34} \\
35 & ANAVI Light pHAT é uma placa add-on Raspberry Pi para controlar tiras de LED RGB de 12 V & \pageref{fig:35} \\
36 & controlador de motor de uso geral projetado com kicad em torno do microcontrolador ATmega328 e driver de motor L298P & \pageref{fig:36} \\
37 & Projetos de energia elétrica EV Inc. Axiom é um controlador de motor de 100kW+ de alta potência e alto desempenho baseado na plataforma VESC & \pageref{fig:37} \\
38 & O Projeto CIAA é um projeto de hardware e software aberto iniciado para alavancar o ambiente industrial argentino & \pageref{fig:38} \\
39 & O Crazyflie 1.0 é uma plataforma de desenvolvimento de voo aberto & \pageref{fig:39} \\
40 & Driverino-Shield é um controlador simples e de baixo consumo para motores BLDC sensorizados & \pageref{fig:40} \\
41 & Glasgow é uma ferramenta para explorar interfaces digitais, destinada a desenvolvedores embarcados, engenharia reversa e outras utilidades & \pageref{fig:41} \\
42 & HackRF One da Great Scott Gadgets é um periférico de Rádio Definido por Software capaz de transmitir ou receber sinais de rádio de 1 MHz a 6 GHz & \pageref{fig:42} \\
43 & CSEduino é a resposta para uma placa DIY Arduino de custo muito baixo & \pageref{fig:43} \\
44 & HADES FCS é um sistema de controle de voo de código aberto para veículos aéreos não tripulados (UAVs) projetado completamente do zero & \pageref{fig:44} \\
45 & O JEMTech-ILDA-Node é um DAC de 16 bits e 6 canais para ILDA de até 400 kpps & \pageref{fig:45} \\
46 & Quadro Suculento-JuicyBoard do plugg.ee Labs é uma plataforma de robótica modular de código aberto que pode ser usada para construir impressoras 3D personalizadas, máquinas CNC e muito mais & \pageref{fig:46} \\
47 & LABDOS-espectrômetro de radiação ionizante & \pageref{fig:47} \\
48 & NUCO-V-placa de desenvolvimento compatível com NUCLEO-64© para as séries STM32F7 e STM32H7 com Black Magic Probe integrado & \pageref{fig:48} \\
49 & Open Smartwatch é um relógio inteligente completamente aberto (ECAD/MCAD e software) & \pageref{fig:49} \\
50 & display POV de código aberto com resolução de 128 pixels e uma taxa de quadros máxima de 20 FPS & \pageref{fig:50} \\
51 & O Picuno é a mistura do RP2040 e do Arduino UNO & \pageref{fig:51} \\
52 & Saturn é uma placa de desenvolvimento FPGA fácil de usar com Xilinx Spartan-6 FPGA & \pageref{fig:52} \\
53 & ScopeFun é uma plataforma de instrumentação tudo-em-um de código aberto & \pageref{fig:53} \\
54 & SmartPrintCoreH7x-controladora 3D & \pageref{fig:54} \\
55 & dosímetros SPACEDOS são espectrômetros e dosímetros de radiação ionizante baseados em detectores semicondutores de código aberto & \pageref{fig:55} \\
56 & O TERES I é um laptop de hardware e software de código aberto & \pageref{fig:56} \\
57 & Tokay Lite: Câmera ESP32 Edge AIMAXLAB.IO & \pageref{fig:57} \\
58 & O primeiro satélite da UPSat Cubesat 2U construído e entregue pela Libre Space & \pageref{fig:58} \\
59 & O arsenal USB 2 da Inverse Path é um projeto de hardware de código aberto que implementa um computador do tamanho de um pen drive & \pageref{fig:59} \\
60 & Label nas trilhas de comunicação USB com diodos de proteção & \pageref{fig:60} \\
61 & Definição dos labels do par diferencial e atribuição a classe DP\_90R & \pageref{fig:61} \\
62 & Label nas trilhas de comunicação USB sem diodos de proteção & \pageref{fig:62} \\
63 & Estrutura com parâmetros para o controle de impedância & \pageref{fig:63} \\
64 & Configuração da placa com edição dos parâmetros conforme compatibilidade da JLCPCB & \pageref{fig:64} \\
65 & Página de compatibilidades da JLCPCB. Link para guia e calculadora de impedância & \pageref{fig:65} \\
66 & Trilhas não revestidas (uncoated) par diferencial & \pageref{fig:66} \\
67 & Estrutura com parâmetros para o controle de impedância & \pageref{fig:67} \\
68 & Constantes dielétricas obtidas na páginas de controle de impedância do JLCPCB & \pageref{fig:68} \\
69 & Parâmetros da classe DP\_90R & \pageref{fig:69} \\
70 & Calculadora de impedância da Sierra Circuits & \pageref{fig:70} \\
71 & Largura da trilha para a classe DP\_90R & \pageref{fig:71} \\
72 & Roteamento do par diferencial & \pageref{fig:72} \\
73 & Comprimento da trilha do par diferencial & \pageref{fig:73} \\
74 & Comprimento da trilha do par diferencial & \pageref{fig:74} \\
75 & Ajuste do comprimento desejado & \pageref{fig:75} \\
76 & Maior trilha com comprimento 33 & \pageref{fig:76} \\
77 & Ajuste de comprimento da trilha & \pageref{fig:77} \\
78 & Trilhas par diferencial com comprimentos iguais e impedância ajustada & \pageref{fig:78} \\
79 & Exemplo esquemático para sinal discreto de campo tipo P & \pageref{fig:79} \\
80 & Exemplo esquemático para sinal CA discreto de campo & \pageref{fig:80} \\
81 & Exemplo esquemático para saída digital CC & \pageref{fig:81} \\
82 & Exemplo esquemático para saída digital com acionamento de carga CA. 97 & \pageref{fig:82} \\
83 & Exemplo esquemático para saída digital com acionamento de carga CA ou CC & \pageref{fig:83} \\
\end{longtable}
\endgroup
